\section{Критерий Арцела-Асколи}
\begin{definition}
    Пусть $(K, \rho)$ -- метрический компакт.
    Тогда $C(K)$ -- пространство непрерывных функций $f\colon K \to \R$.
\end{definition}
\begin{note}
    В силу непрерывности всякой $f \in C(K)$, верно что $f(K)$ -- компакт в $\R$ -- ограниченное и замкнутое множество.
    Значит $x(t)$ достигает минимального и максимального значения:
    \[
        x(t_*) = \max_{K} x, \quad x(t_{**}) = \min_{K} x.
    \]
    И, следовательно $|x|$ достигает на $C(K)$ достигает своего максимального значения.

    Таким образом мы можем ввести $C$-норму как
    \[
        \|x\|_C = \max_{K} |x|.
    \]

    Пространство $C(K)$ с такой нормой является банаховым.
    \begin{proof}
        Пусть $\{f_n\} \subset C(K)$ -- фундаментальная последовательность.
        Заметим, что $\forall x \in K$ последовательность $\{f_n(x)\}$ -- фундаментальна.
        Следовательно, она имеет предел и корректно определена функция $f\colon K \to \R$, что
        \[
            \forall x \in K \hookrightarrow f_n(x) \rightarrow f(x).
        \]
        Заметим, что $f_n$ сходится к $f$ равномерно, в силу равномерной фундаментальности:
        \[
            \forall \epsilon > 0 \exists N\colon \forall n,m \geq N \hookrightarrow \|f_n - f_m\| \leq \epsilon \Rightarrow \|f_n - f\| \leq \epsilon.
        \]
        Поскольку равномерный предел непрерывных -- непрерывен, то $f \in C(K)$ и полнота $C(K)$ показана.
    \end{proof}
\end{note}
\begin{theorem}[Кантор]
    Пусть $x \in C(K)$.
    Тогда $x$ -- равномерно непрерывна.
\end{theorem}
\begin{proof}
    В силу непрерывности $x \in C(K)$
    \[
        \forall \epsilon > 0 \forall t \in K \  \exists \delta(\epsilon, t)\colon \rho(t, \tau) \leq \delta \Rightarrow |x(t) - x(\tau)| < \epsilon.
    \]

    Заметим, что
    \[
        P = \left\{O_{\delta(t, \epsilon)/3}(t) \mid t \in K\right\}
    \]
    является открытым покрытием компакта $K$.

    Тогда существует некоторое конечное подпокрытие $\{O_{\delta(t_i, \epsilon)/3}(t_i)\}_{i = 1}^{N}$.
    Рассмотрим $\delta = \frac{1}{3}\min_{i} \delta(t_i, \epsilon) > 0$.

    Пусть $t, \tau \in K\colon \rho(t, \tau) \leq \delta$.
    Тогда $t \in O_{\delta_j}(t_j)$ и $\tau \in O_{\delta_i}(t_i)$.
    Поскольку
    \[
        \rho(x(t_j), x(t_i)) \leq \frac{\delta(t_j, \epsilon)}{3} + \frac{\delta(t_i, \epsilon)}{3} + \delta \leq \delta(t_i, \epsilon),
    \]
    где мы, без ограничения общности считаем, что $\delta(t_i, \epsilon) \geq \delta(t_j, \epsilon)$.
    А значит
    \[
        |x(t) - x(\tau)| \leq |x(t) - x(t_j)| + |x(t_j) - x(t_i)| + |x(t_i) - x(\tau)| \leq 3\epsilon.
    \]
    И равномерная непрерывность показана.
\end{proof}
\begin{note}
    Заметим теперь некоторые необходимые условия вполне ограниченности в $C(K)$.
    Если $S \subset C(K)$ -- вполне ограничено, то, очевидно, что $S$ -- ограничено в $C(K)$, поскольку лежит в конечном объединении ограниченных множеств.

    Пусть зафиксировано $\epsilon > 0$.
    Из вполне ограниченности следует существование конечной $\epsilon$-сети $\{x_k\}_{k = 1}^{N}$.
    Заметим, что по теореме Кантора каждая из $x_k$ -- равномерно непрерывна с $\delta_k = \delta(x_k, \epsilon)$
    Рассмотрим тогда $\delta = \min \delta_k$.

    Поскольку для любого $x \in S$ существует $x_k$ такое, что $\|x - x_k\|_C \leq \epsilon$, то
    \[
        \forall t, \tau \in K\colon \rho(t, \tau) \leq \delta \hookrightarrow |x(t) - x(\tau)| \leq |x(t) - x_k(t)| + |x_k(t) - x_k(\tau)| + |x_k(\tau) - x(\tau)| \leq 3\epsilon.
    \]
    Это свойство называется равностепенной непрерывностью:
    \[
        \forall \epsilon > 0 \exists \delta(\epsilon) > 0 \forall x \in C(K) \ \forall t, \tau \in K \colon \rho(t, \tau) \leq \delta \Rightarrow |x(t) - x(\tau)| \leq \epsilon.
    \]
\end{note}
\begin{example}
    Пусть $S = \{\sin nt \mid n \in \N\}$.
    Оно не является равностепенно непрерывным, поскольку
    \[
        \exists \epsilon = \sin 1 \forall \delta > 0 \exists t = 0 \wedge \tau = \frac{1}{N}\colon \frac{1}{N} < \delta \wedge \sin Nt\colon |\sin 0 - \sin N \cdot \frac{1}{N}| = \sin 1 \geq \epsilon.
    \]
\end{example}
\begin{theorem}[Арцел, Асколи]
    Пусть $S \subset C(K)$ -- ограниченное и равностепенно непрерывное множество.
    Тогда $S$ -- вполне ограниченное множество.
\end{theorem}
\begin{proof}
    Так как $S$ -- ограничено, то $\exists R > 0$ такое, что
    \[
        \forall x \in S \hookrightarrow \|x\|_C \leq R.
    \]
    И так как оно равностепенно непрерывно, то
    \[
        \forall \epsilon > 0 \exists \delta(\epsilon) > 0 \ \forall x \in C(K) \ \forall t, \tau \in K\colon \rho(t, \tau) \leq \delta \Rightarrow |x(t) - x(\tau)| \leq \epsilon.
    \]

    В компакте $K$ существует конечная $\delta(\epsilon)$-сеть $\{t_j\}_{j = 1}^{N}$.
    Пусть $S_\epsilon \subset \R^N$ состоит из точек
    \[
        S_{\epsilon} = \left\{\begin{pmatrix*} x_(t_1) \\ x(t_2) \\ \ldots \\ x(t_N)
        \end{pmatrix*} \mid x \in S  \right\}.
    \]
    И всё это содержится в кубе $C_R = \{z \in \R^N \mid |z_i| \leq R \forall i = \overline{1, N}\}$.
    Разделим большой кубик $C_R = [-R, R]^N$ на одинаковые кубики с ребром не более чем $\epsilon$.
    Пусть
    \[
        G = \{\square \mid \square \cap S_\epsilon \neq \emptyset\}.
    \]
    Тогда
    \[\forall \square \in G \exists x_{\square} \in S\colon \begin{pmatrix*} x_{\square}(t_1) \\ \ldots \\ x_{\square}(t_k)
    \end{pmatrix*} \in S_\epsilon \cap \square.
    \]
    Рассмотрим множество $\{x_\square \mid \square \in G\} \subset S$ и покажем, что оно является конечной сетью.

    Пусть $x \in S$.
    Тогда $\begin{pmatrix*} x(t_1) \\ x(t_N)
    \end{pmatrix*} \in S_\epsilon \cap \square$ для некоторого $\square \in G$.
    А значит для этого кубика есть $x_\square$.

    Посмотрим, насколько $x$ удалено от $x_\square$.
    Заметим, что так как $\{t_k\}$ -- $\delta(\epsilon)$-есть, то $\forall t \in K \exists t_k\colon \rho(t, t_k) \leq \delta(\epsilon)$.
    Тогда
    \[
        |x(t) - x_\square(t)| \leq |x(t) - x(t_k)| + |x(t_k) - x_\square(t_k)| + |x_{\square}(t_k) - x_{\square}(t)| \leq 3\epsilon.
    \]
    Таким образом
    \[
        \|x(t) - x_{\square}(t)\|_C \leq 3\epsilon
    \]
    И у нас есть конечная $3\epsilon$-сеть -- что и даёт вполне ограниченность $S$.
\end{proof}
\begin{example}
    Пусть у нас есть $K \subset \R^n$ -- выпуклый компакт.
    Рассмотрим множество $S \subset C(K)$, являющееся подмножеством гладких на $K$ функций.

    Пусть $x, y \in K$ и $t \in [0, 1]$.
    Тогда
    \[
        u(y) - u(x) = \int_0^{1} \frac{d}{dt}u(x + t(y - x))dt = \int_0^{1} \langle \nabla u(x + t(y - x)), y - x \rangle dt.
    \]
    А значит:
    \[
        |u(y) - u(x)| \leq \int_{0}^1 |\nabla u(x + t(y - x))|dt \cdot |y - x|.
    \]
    И для вполне ограниченности $S$ нам достаточно равномерной ограниченности градиентов функций из $S$ и ограниченности самого $S$.
\end{example}

\chapter{Теория линейных операторов в нормированных пространствах}
\section{Норма и непрерывность оператора.}
Будем работать с нормированными пространствами $X$, $Y$ и линейными отображениями между ними, называемыми операторами.
\begin{proposition}
    Для линейного оператора $A\colon X \to Y$ следующие условия эквивалентны.
    \begin{enumerate}
        \item $A$ -- непрерывен;
        \item $A$ -- непрерывен в $0$;
        \item $A$ -- ограничен.
    \end{enumerate}
\end{proposition}
\begin{proof}
    Импликация $1) \Rightarrow 2)$ очевидна.

    Пусть $x \in X$ и отображение непрерывно в 0, то есть
    \[
        \forall U(0) \subset Y \ \exists V(0) \subset X\colon AV(0) \subset U(0).
    \]
    Но, тогда, произвольная окрестность $U(Ax) = Ax + U(0)$, а значит $\exists V(0)\colon Ax + AV(0) \subset Ax + U(0) \Rightarrow \exists V(x)\colon AV(x) \subset U(Ax)$, в силу непрерывности линейных операций.
    Что и даёт нам искомую непрерывность в каждой точке.

    Ограниченность оператора $A$ означает, что $AB_1(0)$ -- ограниченное множество.
    Очевидно, что если $A$ -- непрерывен, то $AB_1(0)$ -- ограничен, для $\epsilon = 1$ существует $\delta(1)$ такой, что
    \[
        AO_{\delta(1)}(0) \subset O_2(0) \subset Y.
    \]
    Но $AB_{\delta(1)/2}(0) \subset AO_{\delta(1)}(0)$, что и даёт ограниченность образа единичного шара.

    Пусть теперь образ единичного шара ограничен, то есть
    \[
        AB_1(0) \subset O_R(0) \subset Y.
    \]
    Тогда $\forall \epsilon > 0$
    \[
        \exists AO_{\epsilon/R}(0) \subset AB_{\epsilon/R}(0) \subset O_\epsilon(0),
    \]
    и оператор непрерывен в 0.
\end{proof}
\begin{definition}
    Пусть $A\colon X \to Y$ -- непрерывный оператор.
    Нормой оператора будем называть
    \[
        \|A\| = \sup_{\|x\| \leq 1} \|Ax\|.
    \]
\end{definition}
\begin{note}
    Тогда $A$ -- непрерывен тогда и только тогда когда $\|A\| < +\infty$.

    Очевидно, что
    \[
        \|A\| = \sup_{\|x\| = 1} \|Ax\| = \sup_{x \neq 0} \frac{\|Ax\|}{\|x\|}.
    \]

    Тогда
    \[
        \|Ax\| \leq \|A\|\|x\|.
    \]
\end{note}
\begin{example}
    Пусть $X = Y = L^2(E)$, где $E \subset\R^m$ вместе с мерой Лебега.
    Пусть $K \in L^2(E \times E)$, то есть
    \[
        \int_{E \times E} |K|^2 d\mu \otimes \mu < +\infty.
    \]
    По теореме Фубини
    \[
        \int_{E \times E} |K|^2 d\mu \otimes \mu = \int_{E} \left(\int_{E} |K|^2 d\mu \right)d\mu.
    \]
    А значит функция $\int_{E} |K(x, y)|^2 d\mu(y)$ почти всюду конечна, и, таким образом, $K(x, \cdot) \in L^2(E)$ для почти всех $x \in E$.

    Определим оператор $A\colon L^2(E) \to L^2(E)$ как
    \[
        (Au)(x) = \int_{E} K(x, y) u(y) dy.
    \]
    В силу неравенства Гёльдера для всякой $u$ он почти всюду конечен и, следовательно, корректно задаёт некоторую функцию.
    Она действительно будет лежать в $L^2$ в силу того, что и $K$, и $u$ интегрируемы с квадратом.
    Более того, полученный оператор непрерывен:
    \[
        \|Au\|^2_2 = \int_{E} \left(\int_{E}|K(x, y) u(y)|dy\right)^2 dx \leq \|K\|^2_2 \|u\|_2^2.
    \]
    и его норма не превосходит норму $K$.
\end{example}
\begin{example}[Классический оператор Вольтера]
    Пусть $A\colon L^2([0, 1]) \to L^2([0, 1])$ определяется как
    \[
        (Au)(x) = \int_0^{x} u(y)dy.
    \]
    Это оператор с ядром $K(x, y) = \Ind_{0 \leq y \leq x \leq 1}$, следовательно он является непрерывным и его норма не превосходит $\|K\| \leq \frac{1}{\sqrt{2}}$.

    \begin{notet}
        Дальше идёт точное нахождение нормы оператора при помощи разложения по базису.
        Будет позже.
    \end{notet}
\end{example}

\begin{definition}
    Пространство ограниченных операторов между нормированными пространствами $X$ и $Y$ будем обозначать $\LL(X, Y)$.
    Оно является нормированным вместе с операторной нормой.

    Топология нормы называется <<равномерной операторной топологией>>.
\end{definition}
\begin{proposition}
    Пусть $X$, $Y$ -- нормированные пространства.
    \begin{enumerate}
        \item Если $Y$ -- полно, то $\LL(X, Y)$ -- полно;
        \item Если $Y$ -- неполно, а $X$ нетривиально, то $\LL(X, Y)$ -- неполно.
    \end{enumerate}
\end{proposition}
\begin{proof}
    Пусть $Y$ -- полно.
    Рассмотрим последовательность ограниченных операторов $\{A_n\}$, фундаментальную по операторной норме.
    Тогда $\forall x \in X$ последовательность $\{A_n x\}$ -- фундаментальна (для каждого фиксированного $x$) и, следовательно, сходится к некоторому $Tx$ в силу полноты $y$.
    Таким образом у нас определён линейный оператор
    \[
        Tx = \lim A_n x, \quad \forall x \in X.
    \]
    Покажем его ограниченность.
    В силу неравенства треугольника $|\|A_n\| - \|A_m\|| \leq \|A_n - A_m\|$, а, значит, и последовательность норм операторов является фундаментальной.
    Поэтому она имеет некоторый предел $a$.
    Заметим, что для всякого $x \in B_1(0) \subset X\colon$
    \[
        \|Tx\| = \lim \|A_n x\| \leq \lim\|A_n\|\|x\| = a\|x\|.
    \]
    А значит оператор $T$ -- ограниченный и имеет норму не больше чем $a$.

    Осталось показать, что $A_n$ сходятся к $T$ по операторной норме:
    \[
        \|A_n x - A_m x\| \leq \|A_n - A_m\| \leq \epsilon,
    \]
    с другой стороны, если мы устремим $m \rightarrow \infty\colon$
    \[
        \|A_n x - Tx\| \leq \epsilon.
    \]
    Взятие супремума по $x$ завершает доказательство первого пункта.

    Докажем второй пункт теоремы.
    В силу нетривиальности $X$ существует $x_0 \in X \setminus \{0\}$.
    В силу неполноты $Y$ существует некоторая фундаментальная последовательность $\{y_m\}$, не имеющая предела.
    Рассмотрим $L_0$ -- пространство, натянутое на $x_0$.
    Оно состоит из точек вида $t x_0, t \in \R$.
    Определим $f_0\colon L_0 \to \R, \ tx_0 \mapsto t$.
    Такой функционал, очевидно, непрерывен.
    Мы можем продолжить его на всё $X$ с сохранением нормы при помощи теоремы Хана-Банаха и получить функционал $f\colon X \to \R$.
    Определим $A_n$ как
    \[
        A_n(x) = f(x)y_n.
    \]
    Он является непрерывным и даёт фундаментальную по норме последовательность операторов такую, что в точке $x_0$ она не сходится поточечно -- а значит не может сходиться равномерно в принципе.
\end{proof}
\begin{note}
    В частности для всякого нормированного пространства $X$, сопряжённое $X^* = \LL(X, \R)$ является полным по операторной норме.
\end{note}